



\begin{align*}
  \frac{\partial \beta}{\partial t} (t) &= - \begin{pmatrix}
0 & 0\\
-\overline{\mu} & -\delta
\end{pmatrix}
\begin{pmatrix}
\beta_1(t)\\
\beta_2(t)
\end{pmatrix}
- \begin{pmatrix}
0\\
1
\end{pmatrix}
(\theta(\beta(t)) - 1) \\
\frac{\partial \alpha}{\partial t} (t) &= r- K_0 \begin{pmatrix}
\beta_1(t)\\
\beta_2(t)
\end{pmatrix}
 - \frac{1}{2} \beta_1^2 \sigma^2  
,\end{align*}
where  $\theta(c)$ is the jump transform, defined by:
\[ 
  \theta(c) = \int_{\R^{n}} \ee^{c \cdot z} \d \nu(z),\quad  c \in \C^{n}
.\]
No analytical solution, however, they can be solved fast by numerical routines such as the
Runge-Kutta method.



